  The group \(N\) is then generated by two commuting commutative
  groups, namely \(Z\) and the group generated by \(x\), and is
  therefore commutative.  The group \({}_{\ell^n}N\) is of
  multiplicative type and is characteristic in \(N\); hence it is
  invariant in \(G\), and consequently central because \(G\) is
  connected (\SGARefLink{sga3:ix:corollary:5:5}{Exposé IX, 5.5}).  Thus \({}_{\ell^n}N\) is contained
  in \(Z\), contradicting the fact that its image in \(G'\)
  contains \(\boldsymbol\mu_\ell\).

\item[(ii)] \SGARefTarget{sga3:xvii:lemma:7:3:2:proof:item:ii}
  Suppose that \(\ell=p\).  There is then an integer \(n\) such
  that the image of \(F_nN\) in \(G'\) contains
  \(\boldsymbol\mu_p\), and hence an exact sequence
  \[
    1\longrightarrow K\longrightarrow F_nN
      \longrightarrow\boldsymbol\mu_p\longrightarrow1. \tag{*}\label{sga3:xvii:lemma:7-3-2:equation:star}
  \]
  The group \(K\) is contained in \(Z\), and is therefore
  commutative.  It follows from \SGARefLink{sga3:xvii:theorem:7:2:1}{7.2.1}, \SGARefLink{sga3:xvii:proposition:7:1:1:proof:ii:item:b}{7.1.1(b)}, and \SGARefLink{sga3:xvii:lemma:5:5:1}{5.5.1} that
  there is a subgroup of multiplicative type in \(F_nN\) whose
  image in the quotient by \(K\) is \(\boldsymbol\mu_p\).  As in
  \SGARefLink{sga3:xvii:subsection:6:4:item:i}{(i)}, we deduce that \(F_nN\) is commutative.  Its
  multiplicative-type component (\SGARefLink{sga3:xvii:theorem:7:2:1}{7.2.1}) is characteristic in
  \(F_nN\), hence invariant in \(G\), and therefore central
  (\SGARefLink{sga3:ix:corollary:5:5}{Exposé IX, 5.5}).  Since its image in \(G'\) contains
  \(\boldsymbol\mu_p\), this is a contradiction.
\end{enumerate}

\SGASection{A}{Appendix I. Hochschild Cohomology and Extensions of
Algebraic Groups}
\SGARefTarget{sga3:xvii:appendix:A}

\SGANextSectionTarget{sga3:xvii:subsection:a-1}\SGASubsection{A.1}{Definition of the Cohomology Groups}
\SGARefTarget{sga3:xvii:appendix:A:subsection:1}

Let \(k\) be a field, let \(G\) be an algebraic \(k\)-group, and
let \((\mathrm{Ab})_G\) be the abelian category of commutative
algebraic \(k\)-groups on which \(G\) acts.  If
\(A\in\operatorname{Ob}(\mathrm{Ab})_G\), the functor
\[
  h_A:(\operatorname{Sch}/k)^\circ\longrightarrow
      \operatorname{Ens}
\]
canonically defined by \(A\) is a \(G\)-\(\mathbb Z\)-module in
the sense of \SGARefLink{sga3:i:definition:3-2}{Exposé I, 3.2}.  We may therefore consider the
standard complex \(C^\bullet(G,A)\) of algebraic cochains of \(G\)
with values in \(A\) (\SGARefLink{sga3:i:subsection:5-1}{Exposé I, 5.1}), together with the group
\(Z^i(G,A)\) of \(i\)-cocycles, the group \(B^i(G,A)\) of
\(i\)-coboundaries, and the \(i\)-th cohomology group
\(H^i(G,A)\).  As usual, \(H^0(G,A)\) is identified with the
group \(A^G(k)\), where \(A^G\) is the \(k\)-functor of invariants
of \(A\) under \(G\).  The group \(H^1(G,A)\) classifies the
trivial principal homogeneous spaces under \(A\) on which \(G\)
acts.

In general, the functor
\[
  A\longmapsto H^\bullet(G,A)
\]
from \((\mathrm{Ab})_G\) to \(\mathrm{Ab}\) is not a
cohomological functor.  Nevertheless, we have the following
proposition.

\SGAStatement{\SGAInlineRefTarget{sga3:xvii:proposition:a-1-1}Proposition A.1.1}{Let
\[
  0\longrightarrow A'\xrightarrow{u}A\xrightarrow{v}A''
    \longrightarrow0
\]
be an exact sequence in \((\mathrm{Ab})_G\).  Then:
\begin{enumerate}
\item[(a)] \SGARefTarget{sga3:xvii:proposition:A:1:1:item:a}
  if \(v\) has a section---that is, if there is a \(k\)-morphism
  of preschemes \(s:A''\to A\) such that
  \(v s=1_{A''}\)---then there is the usual exact cohomology
  sequence
  \[
    \cdots\longrightarrow H^i(G,A)\longrightarrow H^i(G,A'')
      \xrightarrow{d}H^{i+1}(G,A')\longrightarrow\cdots;
  \]
\item[(b)] \SGARefTarget{sga3:xvii:proposition:A:1:1:item:b}
  if \(A(k)\to A''(k)\) is surjective, then the sequence
  \[
    0\longrightarrow A'^G(k)\longrightarrow A^G(k)
      \longrightarrow A''^G(k)\xrightarrow{d}H^1(G,A')
      \longrightarrow H^1(G,A)\longrightarrow H^1(G,A'')
      \tag{1}\label{sga3:xvii:proposition:a-1-1:equation:1}
  \]
  is exact.
\end{enumerate}}
\SGARefTarget{sga3:xvii:proposition:A:1:1}

\SGAProof{Proof}
For \SGARefLink{sga3:xvii:appendix:A}{(a)}, observe that the existence of a section implies the
exactness of the sequence of complexes
\[
  0\longrightarrow C^\bullet(G,A')\longrightarrow C^\bullet(G,A)
    \longrightarrow C^\bullet(G,A'')\longrightarrow0.
\]
